% ============================================================
%  anexos.tex — Anexos del documento
% ============================================================

\chapter{Encuesta de Perspectiva sobre la Normalización de Redes Eléctricas en un Contexto Técnico y Social}
\label{anx:encuesta}

La encuesta se realizó mediante un formulario digital asincrónico (Google Forms) con el fin de facilitar la recolección de información, considerando la disponibilidad de los profesionales y su residencia en distintas ciudades del territorio nacional. El objetivo fue identificar la efectividad y el impacto en los OR y en los usuarios de la implementación de proyectos de normalización, y recopilar, a partir de las experiencias de los participantes, información sobre los cuellos de botella que fundamentan la necesidad de acciones complementarias en la intervención en BS.

% ── Ficha técnica ──────────────────────────────────────────

\section*{Ficha Técnica del Instrumento}
\addcontentsline{toc}{section}{A1. Ficha técnica del instrumento}

\begin{table}[H]
\centering
\caption{Ficha técnica del instrumento aplicado a expertos con experiencia en el sector eléctrico}
\label{tab:ficha_tecnica}
\begin{tabularx}{\textwidth}{>{\bfseries\RaggedRight}p{4cm} >{\RaggedRight\arraybackslash}X}
\toprule
\textbf{Atributo} & \textbf{Detalle} \\
\midrule
\rowcolor{grisclaro}
Tipo de instrumento & Encuesta con preguntas de escala Likert y de respuesta abierta, de respuesta obligatoria. \\
Medio de aplicación & Formulario asincrónico en línea (Google Forms). \\
\rowcolor{grisclaro}
Cantidad de preguntas & 17 preguntas, identificadas como P1 a P17 en el desarrollo del texto. \\
Población objetivo & Profesionales en ingeniería eléctrica, ingeniería civil, ciencias sociales y administración pública con experiencia en proyectos de normalización de redes eléctricas en Colombia. \\
\rowcolor{grisclaro}
Tamaño de la muestra & 20 expertos, identificados como E1 a E20 en el análisis de resultados. \\
Año de aplicación & 2025. \\
\rowcolor{grisclaro}
Tratamiento de datos & Anonimato garantizado, dado que el instrumento se aplica con fines estrictamente académicos. \\
\bottomrule
\multicolumn{2}{l}{\footnotesize\textit{Fuente: Elaboración propia.}}
\end{tabularx}
\end{table}

% ── Perfil de expertos ─────────────────────────────────────

\section*{Caracterización de la Experiencia de los Expertos Encuestados}
\addcontentsline{toc}{section}{A2. Caracterización de los expertos}

\begin{longtable}{>{\centering}p{1cm}
                  >{\RaggedRight}p{2.4cm}
                  >{\centering}p{1.5cm}
                  >{\RaggedRight}p{3.0cm}
                  >{\RaggedRight}p{2.8cm}
                  >{\RaggedRight\arraybackslash}p{2.5cm}}
\caption{Caracterización de la experiencia de los expertos encuestados}
\label{tab:perfil_expertos} \\
\toprule
\textbf{No.} & \textbf{Área de formación [P2]} & \textbf{Años de exp. [P3]} & \textbf{Área de desempeño [P4]} & \textbf{Región [P5]} & \textbf{Tipo de entidad [P6]} \\
\midrule
\endfirsthead
\multicolumn{6}{c}{\small(Continuación Tabla \ref{tab:perfil_expertos})} \\
\toprule
\textbf{No.} & \textbf{Área de formación [P2]} & \textbf{Años de exp. [P3]} & \textbf{Área de desempeño [P4]} & \textbf{Región [P5]} & \textbf{Tipo de entidad [P6]} \\
\midrule
\endhead
\bottomrule
\endfoot
\bottomrule
\endlastfoot

E1  & Administrador Público     & 1--2 & Técnica/Ingeniería          & Caribe                         & Público \\
\rowcolor{grisclaro}
E2  & Administrador Público     & 3--5 & Jurídica/Administrativa     & Caribe                         & Público \\
E3  & Ciencias sociales y humanas & 1--2 & Social/Comunitaria        & Caribe, Andina                 & Público \\
\rowcolor{grisclaro}
E4  & Ciencias sociales y humanas & 3--5 & Social/Comunitaria        & Caribe, Pacífico               & Público \\
E5  & Ingeniero civil           & 6--8 & Técnica/Ingeniería          & Pacífico                       & Público y privado (mayor exp. privado) \\
\rowcolor{grisclaro}
E6  & Ingeniero Eléctrico       & 1--2 & Jurídica/Administrativa     & Caribe, Pacífico, Andina, Orinoquía, Amazonía & Público \\
E7  & Ingeniero Eléctrico       & 1--2 & Técnica/Ingeniería          & Caribe, Andina                 & Público \\
\rowcolor{grisclaro}
E8  & Ingeniero Eléctrico       & 1--2 & Técnica/Ingeniería          & Caribe, Pacífico               & Público \\
E9  & Ingeniero Eléctrico       & 1--2 & Técnica/Ingeniería          & Caribe, Andina                 & Público y privado (mayor exp. privado) \\
\rowcolor{grisclaro}
E10 & Ingeniero Eléctrico       & 1--2 & Técnica/Ingeniería          & Caribe, Pacífico, Andina, Orinoquía & Público \\
E11 & Ingeniero Eléctrico       & 3--5 & Técnica/Ingeniería          & Caribe, Pacífico, Amazonas     & Público \\
\rowcolor{grisclaro}
E12 & Ingeniero Eléctrico       & 3--5 & Comercial/Pérdidas de energía & Caribe, Andina               & Público y privado (mayor exp. privado) \\
E13 & Ingeniero Eléctrico       & 3--5 & Técnica/Ingeniería          & Andina                         & Público y privado (mayor exp. público) \\
\rowcolor{grisclaro}
E14 & Ingeniero Eléctrico       & 6--8 & Técnica/Ingeniería          & Caribe, Orinoquía              & Público \\
E15 & Ingeniero Eléctrico       & 6--8 & Técnica/Ingeniería          & Caribe                         & Público \\
\rowcolor{grisclaro}
E16 & Ingeniero Eléctrico       & 6--8 & Técnica/Ingeniería          & Caribe, Pacífico               & Público \\
E17 & Ingeniero Eléctrico       & $>$9 & Comercial/Pérdidas de energía & Andina                       & Público y privado (mayor exp. público) \\
\rowcolor{grisclaro}
E18 & Ingeniero Eléctrico       & $>$9 & Técnica/Ingeniería          & Caribe, Andina                 & Público y privado (mayor exp. público) \\
E19 & Ingeniero Eléctrico       & $>$9 & Técnica/Ingeniería          & Caribe, Andina                 & Público y privado (mayor exp. público) \\
\rowcolor{grisclaro}
E20 & Ingeniero Eléctrico       & $>$9 & Técnica/Ingeniería          & Caribe                         & Público y privado (mayor exp. público) \\
\end{longtable}
\textit{Fuente: Elaboración propia.}

% ── Tabla 7: Impacto en la red ─────────────────────────────

\section*{Resultados: Impacto en la Red Eléctrica y Percepción de los Usuarios}
\addcontentsline{toc}{section}{A3. Resultados -- impacto en la red y percepción}

\begin{longtable}{>{\centering}p{0.7cm}
                  >{\centering}p{1.1cm}
                  >{\centering}p{1.1cm}
                  >{\RaggedRight}p{5.2cm}
                  >{\RaggedRight\arraybackslash}p{5.2cm}}
\caption{Importancia e impacto en la red eléctrica y en la prestación del servicio de energía de los proyectos de normalización}
\label{tab:encuesta_red} \\
\toprule
\textbf{No.} &
\textbf{Reducción PT [P7]} &
\textbf{Disminución fallas [P9]} &
\textbf{¿Considera importante la normalización para reducir accidentes eléctricos? Explique [P10]} &
\textbf{¿Cómo perciben los usuarios el cambio en la calidad y confiabilidad del servicio? [P11]} \\
\midrule
\endfirsthead
\multicolumn{5}{c}{\small(Continuación Tabla \ref{tab:encuesta_red})} \\
\toprule
\textbf{No.} & \textbf{P7} & \textbf{P9} & \textbf{P10} & \textbf{P11} \\
\midrule
\endhead
\bottomrule
\endfoot
\bottomrule
\endlastfoot

E1  & 2 & No se conoce & Sí, para evitar incendios o riesgos eléctricos. & Agradecen tener el servicio de energía pero no quieren pagar por él. \\
\rowcolor{grisclaro}
E2  & 2 & No se conoce & Sí, porque hay una gestión e intervención adecuada de las redes eléctricas. & Satisfacción asociada a la estabilidad y confiabilidad del servicio público. \\
E3  & 2 & No se conoce & Sí. La mejora de las condiciones físicas reduce las posibilidades del uso informal del servicio. Con una socialización y capacitación adecuadas, se puede sensibilizar frente a los riesgos a la vida. & Lo perciben inicialmente a través de la reducción o nulos cortes no programados y el no daño de sus electrodomésticos. Sobre la confiabilidad, las comunidades suelen desconocer su significado técnico y lo asocian con la veracidad en el cobro. \\
\rowcolor{grisclaro}
E4  & 3 & Sí & Sí, cuando se da el manejo de redes por personal no cualificado, suelen ocurrir incidentes que ponen en riesgo a la comunidad. & Usualmente es positivo, en cuanto la normalización supere la media y se den las inversiones correspondientes. Si no agrupa a la mayoría, los usuarios manifiestan malestar porque la calidad no mejora considerablemente. \\
E5  & 5 & Sí & Sí, porque la normalización reemplaza instalaciones informales o en mal estado por infraestructura técnica segura conforme al RETIE, reduciendo riesgos de cortocircuitos, sobrecargas, electrocución e incendios. & En general perciben mejora: mayor continuidad, reducción de interrupciones, estabilidad en el voltaje. Sin embargo, el cambio puede requerir adaptación porque no siempre proviene de comunidades con cultura de pago. \\
\rowcolor{grisclaro}
E6  & 3 & No se conoce & Sí. & No sé. \\
E7  & 4 & No se conoce & Sí, ya que al normalizar las redes se busca transformar las derivaciones ilícitas en conexiones que cumplan los criterios normativos de seguridad, reduciendo la posibilidad de incidentes eléctricos. & Sin información. \\
\rowcolor{grisclaro}
E8  & 4 & No se conoce & Sí. El hecho de tener una red normalizada disminuye el riesgo de instalaciones sin protecciones ni manejo correcto de cargas que puedan afectar directamente a los usuarios. & No poseo información de primera mano. \\
E9  & 4 & No se conoce & Sí, la consideraría muy importante. Una red normalizada es más fácil de inspeccionar, mantener y ampliar sin improvisación, lo que baja el riesgo de accidentes a largo plazo. & Normalmente perciben la mejora cuando disminuyen los cortes y las fluctuaciones que afectan sus electrodomésticos. Pero también puede haber percepción negativa si los eventos siguen siendo frecuentes, incluso si hay mejoras técnicas no evidentes para el usuario. \\
\rowcolor{grisclaro}
E10 & 5 & No se conoce & Sí, ya que la normalización garantiza que las instalaciones cumplan con los estándares técnicos del RETIE. & Los usuarios perciben la normalización positivamente: mayor continuidad en el servicio y menos interrupciones. \\
E11 & 2 & No & Sí, aunque debe considerarse la cantidad de usuarios normalizados que reinciden en la subnormalidad con el fin de no pagar. Las condiciones de normalidad aportan mayor seguridad, pero no garantizan que así se mantengan las redes. & Creen que lo perciben con un servicio más confiable y mejor funcionamiento de sus electrodomésticos. \\
\rowcolor{grisclaro}
E12 & 3 & No & Considero que estos barrios, dependiendo del número de habitantes, no representan una carga significativa para el sistema eléctrico por sus bajos consumos. Sí se identificaron riesgos importantes entre quienes se conectaban fraudulentamente. & En muchos casos se les cobra entre \$400.000 y más por medidores, sellos y costos de conexión. Si a esto se suma el CNR, muchas familias prefieren no conectarse porque no tienen recursos para asumir esos costos iniciales. \\
E13 & 4 & Sí & Sí, ya que medidas comunitarias pueden conllevar malas conexiones y conexiones fraudulentas que generen siniestros. & De mejor manera, ya que la acometida es específica para cada usuario. \\
\rowcolor{grisclaro}
E14 & 1 & Sí & Sí, aunque se deberían revisar los datos de la Superservicios. & La calidad y confiabilidad aumenta a pesar de la falta de percepción por parte del usuario. \\
E15 & 4 & No se conoce & Sí, debemos adecuarnos al reglamento y las normas técnicas. & Considero que hay un problema social o una cultura del no pago, y allí radican los problemas de los barrios subnormales. \\
\rowcolor{grisclaro}
E16 & 5 & Sí & Sí, la normalización es clave para reducir diferentes tipos de riesgos asociados a la prestación del servicio, ya que elimina conexiones inseguras, garantiza el cumplimiento del RETIE y mejora la seguridad de las instalaciones residenciales. & Los usuarios perciben un cambio positivo en términos de menos cortes del servicio, menor daño en sus electrodomésticos y mayor seguridad para las personas. \\
E17 & 2 & Sí & Es absolutamente importante y necesaria; es la razón principal de la normalización. Desde el punto de vista comercial, si se factura una medida comunitaria, se está contemplando el nivel de pérdidas y el problema se vuelve de cartera. & Lo consideran un derecho y una obligación del Estado; desde el punto de vista de calidad de vida, lo consideran una mejora importante. \\
\rowcolor{grisclaro}
E18 & 4 & Sí & Sí, la normalización se realiza cumpliendo la normatividad vigente, en especial el RETIE, cuyo objeto es minimizar o eliminar los riesgos eléctricos. & Los usuarios perciben un mejor servicio en general. \\
E19 & 4 & Sí & Sí, porque se reducen las malas prácticas. & --- \\
\rowcolor{grisclaro}
E20 & 5 & Sí & Sí, la normalización permite que las redes den cumplimiento a la normativa del RETIE. & Los usuarios muy felices con la confiabilidad del servicio; no tienen interrupciones y consideran que el barrio tiene ahora muchos transformadores, lo que les da tranquilidad. \\
\end{longtable}
\noindent\textit{Fuente: Elaboración propia.}

% ── Tabla 8: Recaudo y facturación ────────────────────────

\section*{Resultados: Impacto en el Recaudo y Perspectiva del Usuario frente a la Facturación Individual}
\addcontentsline{toc}{section}{A4. Resultados -- recaudo y facturación individual}

\begin{longtable}{>{\centering}p{0.7cm}
                  >{\RaggedRight}p{4.5cm}
                  >{\RaggedRight}p{4.5cm}
                  >{\RaggedRight\arraybackslash}p{4.5cm}}
\caption{Impacto en el recaudo y perspectiva del usuario frente al medidor y facturación individual cuando se implementan proyectos de normalización}
\label{tab:encuesta_recaudo} \\
\toprule
\textbf{No.} &
\textbf{¿Considera una mejora en la eficiencia del recaudo al pasar de esquemas comunitarios a individuales? [P8]} &
\textbf{¿Cuál es la posición del usuario frente a ser legalizado (medidor y factura individual)? [P12]} &
\textbf{¿Qué tan crítico es el impacto del cambio en el valor de la factura para un usuario que antes pagaba cuota fija o no pagaba? [P13]} \\
\midrule
\endfirsthead
\multicolumn{4}{c}{\small(Continuación Tabla \ref{tab:encuesta_recaudo})} \\
\toprule
\textbf{No.} & \textbf{P8} & \textbf{P12} & \textbf{P13} \\
\midrule
\endhead
\bottomrule
\endfoot
\bottomrule
\endlastfoot

E1  & No. La cultura del no pago se mantiene aun cuando con el recaudo se pudiera garantizar la prestación del servicio. & No pagar el servicio, puesto que algunos no cuentan con los recursos financieros. & Es crítico, pues esperan que el Estado haga la prestación del servicio por ser un derecho sin que se pague por ello. \\
\rowcolor{grisclaro}
E2  & Sí. & Rechazo asociado al pago del servicio público de energía. & Depende de la variación en la factura y la capacidad de ingresos y pago del hogar. \\
E3  & Sí, parcialmente mejora la percepción de cobro ajustado a la realidad individual. & Mayoritariamente es una posición poco receptiva. Pasa por la deuda histórica que tiene el Estado, además de la relación cultural que han tenido algunas regiones con la electricidad y el sector en general. La precarización de las condiciones de vida de algunos sectores hace que las personas busquen sobrevivir en todas las dimensiones de su vida y vean en la subnormalidad la posibilidad de tener gratuidad en este servicio. & Podría resultar bastante crítico, en sentido de deterioro o negativo, si corresponde a sectores que presentan falencias o imposibilidad de abastecer las demás necesidades básicas. En otros casos, donde se comprueba capacidad de pago, podría resultar de poco impacto crítico. \\
\rowcolor{grisclaro}
E4  & Es relativo a las comunidades, al proceso pedagógico que se adelante y a la claridad en la información de la factura que se entregue. Hay diferentes variables a considerar. & No hay una única posición; varía según el usuario, la comunidad y las acciones de pedagogía y socialización que se adelanten. & Si no se construye una cultura de pago y se adoptan medidas de racionalidad y eficiencia energética, el impacto se siente de manera considerable. Si el OR cobra deudas antiguas al usuario, también se genera un incremento considerable en la factura. \\
E5  & Sí, porque la individualización permite mayor control del consumo, mejora la transparencia en la facturación, reduce pérdidas y fortalece la cultura de pago. No obstante, requiere acompañamiento social para garantizar la adaptación. & Los usuarios muestran una percepción favorable frente a la legalización por las mejoras en calidad, seguridad y confiabilidad. Sin embargo, algunos manifiestan preocupación por los costos del pago individual, por lo que el proceso requiere acompañamiento pedagógico. & Es un impacto alto, porque el usuario asume el costo real del consumo, lo que puede generar afectación económica y resistencia inicial, aunque promueve el uso eficiente de la energía. \\
\rowcolor{grisclaro}
E6  & Sí, porque se pueden identificar la cantidad de unidades residenciales y usuarios individuales que reciben el servicio de energía. & Es necesario que los usuarios cuenten con un servicio confiable y seguro, que evite riesgos a la vida humana frente a accidentes eléctricos. Sin embargo, el desconocimiento, la falta de información y los imaginarios colectivos contribuyen a que la gente tenga una percepción negativa frente a la legalización del servicio. & No sé. \\
E7  & Desde un punto de vista de usuario, es más eficiente el recaudo individual, ya que se tiene pleno conocimiento de que la facturación es únicamente el consumo propio. Sin embargo, para usuarios malintencionados esta práctica es contraproducente, ya que estarían pagando por un servicio que antes era ``gratis''. & En muchas ocasiones, la cultura del no pago está tan arraigada que el hecho de instalar un medidor representa un problema para ellos. Sin embargo, para usuarios afectados por la mala calidad de la red, la nueva infraestructura representa un beneficio sin precedentes. & En general, los proyectos de normalización se dan en barrios de baja estratificación, donde cualquier gasto adicional puede representar un gran esfuerzo. El mayor impacto siempre será para aquellos usuarios conectados de manera ilícita, ya que el cambio será pasar de no pagar a pagar una factura. \\
\rowcolor{grisclaro}
E8  & Sí, pasar a esquemas individuales permite mayor control y trazabilidad para el recaudo. & Supongo que ven dificultades por la implicación de tener que pagar el servicio de energía eléctrica, sumado a la falta de confianza en el OR para zonas con fallas constantes. & Es un impacto considerable, pues los usuarios usualmente son de bajos ingresos y presentan dificultad para el pago del servicio. \\
E9  & Sí, normalmente mejora la eficiencia de facturación y recaudo porque el consumo queda medido y asignado por usuario, reduciendo la ambigüedad de ``quién debe cuánto''. Sin embargo, la mejora en recaudo está de la mano de la capacidad de pago; si esta es baja o hay rechazo social al cobro, puede aumentar la mora aunque la facturación sea más precisa. & En principio creo que estaría a favor, porque un medidor y una factura individual dan claridad sobre el consumo, formalizan el servicio y facilitan la atención del OR. Sin embargo, su aceptación dependería de que el costo de la energía sea manejable y de que realmente mejore la calidad y continuidad del servicio: si no mejora, ¿para qué se legalizan? & Considero que el impacto sería alto, porque pasar de una cuota fija o de no pagar a una factura variable puede desajustar el presupuesto de los usuarios. Además, genera incertidumbre, porque son usuarios que no tienen hábitos de control del consumo. \\
\rowcolor{grisclaro}
E10 & Sí, pasar de una facturación comunitaria a una individual contribuye a una mejora en el recaudo, ya que permite una asignación directa al consumo de cada usuario, promueve mayor responsabilidad y mejora la transparencia. & En algunos casos la percepción puede ser negativa, al tener que pagar facturas que antes no pagaban y tener que moderar el consumo de la energía. & El cambio puede impactar significativamente la economía del hogar, pues pasar de no pagar o pagar una cuota muy baja a una facturación del consumo real puede generar inicialmente resistencia y percepción negativa, así el servicio mejore. \\
E11 & No necesariamente, dado que normalmente la normalización genera que se pase de la subnormalidad a una situación de cartera en normalidad. & Lo que he podido notar es que en la mayoría de casos se presenta resistencia al pago, lo cual genera la acumulación de cartera. & Es realmente crítico considerando que los consumos en las zonas subnormales son elevados, cubiertos solo hasta el consumo de subsistencia por los subsidios, y que una vez normalizados pierden en algunos casos el beneficio FOES. Si la normalización no viene acompañada de concientización y revisión de las instalaciones y hábitos de consumo, el impacto puede ser tan crítico que genere un retorno a la subnormalidad. \\
\rowcolor{grisclaro}
E12 & Considero que, más que representar un aporte económico para las empresas, este proceso tiene un impacto principalmente social. La normalización reduce de manera significativa los riesgos de accidentes y, en muchos casos, se evidenció que algunas personas se apropiaban de las redes y cobraban por realizar conexiones ilegales, operando como ``criminales eléctricos''. El porcentaje de recaudo es mínimo dado que aproximadamente el 50\% de la factura es subsidiado. Se trata de un proceso de dignificación y garantía de derechos, más que de eficiencia financiera. & En muchos casos lo ven de forma negativa y se rehúsan a ser conectados. & Depende del sector y de los cobros asociados a la legalización, pero el impacto es grande. \\
E13 & Sí, ya que se controla de manera más específica el consumo uno a uno y las pérdidas que lleguen a generar. & Algunos se resisten, pero desde que se tenga confiabilidad en el servicio y una buena campaña de concientización, se legalizan. & Es un impacto grande, ya que algunos usuarios en medidas comunitarias no cuentan con los suficientes recursos para pagar una cuenta completa, pero les genera un hábito de pago cuando ya cuentan con ella. \\
\rowcolor{grisclaro}
E14 & Sí, la responsabilidad no cae sobre una persona o representante y se manifiesta mayor recaudo. & Negativa. & N/A. \\
E15 & No, considero que el tema es de carácter social. & Gran parte de los usuarios se oponen por no estar de acuerdo con pagar un servicio. & Afecta el bolsillo, por lo tanto es crítico, pero se debe concientizar a los usuarios. \\
\rowcolor{grisclaro}
E16 & --- & La posición inicial de los usuarios corresponde a cierto temor por un incremento del valor a pagar, desconfianza en el sistema de medición y preocupación por la capacidad de pago. & El impacto es muy crítico ya que puede generar un choque económico al momento de generar una responsabilidad individual económicamente, por lo cual en ciertos usuarios hay una relación mayoritaria con tener que pagar más y no los beneficios de seguridad y confiabilidad. \\
E17 & En las experiencias tenidas se presentan dificultades relacionadas con seguridad y el impedimento de la comunidad de permitir medida telegestionada. Sin embargo, considero que es la mejor opción para las partes, siempre que los procesos sean acompañados de gestión comunitaria. & Rechazo, dado que implica una nueva responsabilidad personal. La medida comunitaria ``es de todos y es de nadie'' y desde el punto de vista legal, no hay a quién cobrarle judicialmente. & Si el usuario paga la medida comunitaria, la individualización es un beneficio para quienes tienen bajo consumo, ya que se ahorran los valores adicionales que pagan por pérdidas internas de la medida comunitaria. Si no pagan la comunitaria, el impacto es total: de cero pesos a un valor, así sea bajo. \\
\rowcolor{grisclaro}
E18 & Sí, acorde con lo establecido en la Ley 142 de 1994, artículo 146: ``La empresa y el suscriptor o usuario tienen derecho a que los consumos se midan; a que se empleen para ello los instrumentos de medida que la técnica haya hecho disponibles; y a que el consumo sea el elemento principal del precio que se cobre''. Así, se crea en los usuarios una mayor cultura de pago al pagar por los servicios realmente recibidos. & En las comunidades más organizadas, se ha evidenciado la solicitud de normalización del servicio. & El valor de la factura puede presentar generalmente reducción cuando se trata de usuarios netamente residenciales. Ahora bien, el impacto se genera al momento de realizar los pagos, ya que si un usuario está acostumbrado al no pago del servicio, cualquier valor le sea cobrado lo verá como un incremento. \\
E19 & Sí porque se mide y se cobra con exactitud. & Inicialmente es difícil, ellos no quieren que se instale el medidor. & Total. Es pasar de cero pesos a un valor, así sea bajo. \\
\rowcolor{grisclaro}
E20 & No, la eficiencia en el recaudo depende de las estrategias que realice el OR, y una puede ser condonando deudas anteriores a la normalización. & El usuario legalizado no le gusta que le coloquen un medidor de energía que le permita efectivamente pagar lo que consume, y más si tiene algún tipo de actividad comercial dentro del barrio normalizado. & Es crítico considerando que algunos usuarios al considerar que no pagaban el servicio no tienen la cultura del ahorro de energía; entonces, después de pagar un valor fijo que puede ser entre \$40.000 a \$50.000, ahora van a pagar \$100.000--\$200.000. \\
\end{longtable}
\noindent\textit{Fuente: Elaboración propia.}

% ── Tabla 9: Estrategias ──────────────────────────────────

\section*{Resultados: Estrategias para Mejorar la Implementación de Proyectos de Normalización}
\addcontentsline{toc}{section}{A5. Resultados -- estrategias y casos documentados}

Las alternativas complementarias referenciadas en la columna [P14] corresponden a las siguientes opciones del instrumento: \textbf{1.} Socialización previa con información clara sobre el proceso y sus beneficios. \textbf{2.} Acompañamiento durante la ejecución de las obras. \textbf{3.} Condonación o refinanciación de deudas históricas previas a la normalización. \textbf{4.} Seguimiento post-normalización con evaluaciones de recaudo y satisfacción. \textbf{5.} Articulación con entidades territoriales y líderes comunitarios.

\begin{longtable}{>{\centering}p{0.7cm}
                  >{\centering}p{2cm}
                  >{\centering}p{1.2cm}
                  >{\RaggedRight}p{3.5cm}
                  >{\RaggedRight\arraybackslash}p{5.0cm}}
\caption{Estrategias para mejorar la implementación de proyectos de normalización y resultado en la ejecución de proyectos de normalización}
\label{tab:encuesta_estrategias} \\
\toprule
\textbf{No.} &
\textbf{Alternativas complementarias [P14]} &
\textbf{Importancia acomp. social [P15]} &
\textbf{Principal obstáculo para que el usuario acepte la normalización [P16]} &
\textbf{Caso donde la normalización tuvo impacto positivo o fracaso. Factor determinante [P17]} \\
\midrule
\endfirsthead
\multicolumn{5}{c}{\small(Continuación Tabla \ref{tab:encuesta_estrategias})} \\
\toprule
\textbf{No.} & \textbf{P14} & \textbf{P15} & \textbf{P16} & \textbf{P17} \\
\midrule
\endhead
\bottomrule
\endfoot
\bottomrule
\endlastfoot

E1  & 1, 2 y 4 & 5 & Falta de información sobre la operación del sistema y desconfianza en el OR. & Los proyectos PRONE implementados en la costa Caribe tuvieron un impacto negativo. Algunas de las causas fueron la desinformación que tenían los usuarios y las demoras en la adecuada ejecución del proyecto. \\
\rowcolor{grisclaro}
E2  & 1, 2 y 4 & 5 & Pago de la factura o incremento en la misma. & No tengo casos de estudio específicos. \\
E3  & Ninguno & 5 & Desconfianza institucional y poco seguimiento por parte de las entidades. & Algunos proyectos en el Caribe. La articulación con las comunidades resulta no contemplar aspectos esenciales para el éxito de los proyectos. La normalización pasa por un cambio de relación entre las personas y la energía como servicio para que sea exitoso. \\
\rowcolor{grisclaro}
E4  & 1, 2, 3 y 5 & 5 & Comprender los derechos y deberes que tengo como usuario. & Patia -- Cauca. Impacto negativo -- poca apropiación por parte de las comunidades. \\
E5  & 1, 2, 3 y 5 & 5 & Creer en el aumento de la factura. & En el departamento de Nariño, proyectos liderados por CEDENAR evidenciaron un impacto positivo en la reducción de pérdidas, mejora en la calidad y continuidad del servicio, y disminución de riesgos eléctricos. En algunos sectores se presentó resistencia inicial por el incremento en el valor de la factura tras la individualización. El factor determinante del éxito ha sido el acompañamiento social y la sensibilización a la comunidad durante la implementación. \\
\rowcolor{grisclaro}
E6  & 1, 2 y 4 & 5 & Falta de información y percepción de aumento de la factura. & --- \\
E7  & 1 y 5 & 5 & Creencia en el aumento de la factura y la desinformación con respecto al objetivo del proyecto. & Sin información. \\
\rowcolor{grisclaro}
E8  & 1 y 2 & 5 & Preocupación por el aumento en el costo de la factura y desconocimiento de la importancia de la normalización desde el punto de vista de seguridad eléctrica. & No conozco casos particulares. \\
E9  & 1 y 5 & 5 & El principal obstáculo debe ser el temor a que aumente el valor de la factura, sumado a la desconfianza en el OR por experiencias previas propias o de terceros. También pesa mucho la falta de información sobre qué implica la normalización. & En 2018 en Isla Fuerte se implementó un sistema híbrido diésel/solar para toda la comunidad. En general el impacto fue positivo: la mayoría de habitantes reportó mejora en su calidad de vida y los comercios pudieron ofrecer servicios de hospedaje más competitivos. No obstante, un grupo de usuarios manifestó incapacidad de pago y una expectativa de recibir la energía de forma gratuita, argumentando que al haber sido conectados por el Estado, era responsabilidad del mismo asumir el suministro. El factor determinante fue la brecha entre el beneficio técnico-social del servicio y la aceptación del esquema de pago, influida en gran medida por la capacidad de pago, la percepción de derecho y la claridad previa sobre reglas tarifarias. \\
\rowcolor{grisclaro}
E10 & 1, 4 y 5 & 5 & El usuario puede tener desconfianza hacia el OR, puede ver aumentos en su factura y no tener información clara sobre cómo se calcula el consumo de energía. & Desde el área en que trabajo evalúo proyectos de normalización que se presentan para ser financiados con recursos públicos. En la práctica se evidencia que lo más importante es llevar la información clara a las comunidades beneficiarias: explicar lo que conlleva la normalización, las fórmulas tarifarias y las mejoras en la seguridad y el servicio que la normalización trae consigo. \\
E11 & 1, 2, 3, 4 y 5 & 5 & Resistencia a asumir el pago de un servicio que no es bien utilizado, además de la desinformación que a veces se genera sobre la gratuidad en el servicio. & No conozco. \\
\rowcolor{grisclaro}
E12 & 2, 3 y 4 & 5 & Desconfianza en el OR, se sienten vulnerados y amenazados. & No tengo como tal un estudio, solo mi experiencia específica. Para los barrios normalizados a largo plazo es beneficioso ya que gracias a esto pueden acceder a educación, ya que cuentan con su factura eléctrica; adicionalmente, otras entidades como acueducto y gas también comienzan procesos de normalización, pueden registrarse al Sisbén, entre otros. \\
E13 & 1 y 4 & 5 & Falta de información e impacto de aumento en el costo de la factura. & En la empresa donde laboro, Electrohuila, se hacen campañas de legalización mejorando las pérdidas de energía y el recaudo de la empresa. \\
\rowcolor{grisclaro}
E14 & 1, 2 y 4 & 5 & La falta de cultura de pago. & --- \\
E15 & 1, 2, 3, 4 y 5 & 5 & Cultura de no pago. & Conozco muchos barrios de la región Caribe donde se debieron cancelar los proyectos por oposición de la comunidad a la medida. \\
\rowcolor{grisclaro}
E16 & 1, 2, 3, 4 y 5 & 5 & Que se le instalará un medidor de energía para facturar el consumo de cada usuario. & Normalmente, los proyectos de normalización de redes eléctricas son bien aceptados en diferentes territorios del departamento del Meta. Hay algunas afectaciones con respecto a la aceptación del proyecto debido al tipo de medición instalado, lo cual ocurre mayoritariamente en la zona Caribe de Colombia. \\
E17 & 1, 3, 4 y 5 & 5 & Asumir responsabilidad personal y directa por su consumo. & Normalización de Brisas del Venado. La comunidad no permitió la instalación de medida AMI. \\
\rowcolor{grisclaro}
E18 & 1, 2, 4 y 5 & 5 & La desinformación generada por los suscriptores que saben que el cobro individual del servicio puede generar mayor valor a pagar por su parte (comercios, pequeñas o medianas industrias y en general aquellos que no propenden por el URE). & Los estudios previos y posteriores son adelantados directamente por los comercializadores de energía que han adelantado estos proyectos. \\
E19 & 4 & 5 & Falta de cultura energética. & --- \\
\rowcolor{grisclaro}
E20 & 1, 2 y 3 & 5 & El tipo de medida que se instale, dado que no están de acuerdo con la medición inteligente. & Un caso de éxito en la normalización fue un proyecto en Caquetá, y un fracaso han sido algunos proyectos de la costa Caribe. El factor determinante fue el esquema de medición. \\
\end{longtable}
\noindent\textit{Fuente: Elaboración propia.}


% ══════════════════════════════════════════════════════════════
\chapter{Municipios de la Región Caribe por Piso Térmico}
\label{anx:municipios_clima}
% ══════════════════════════════════════════════════════════════

Los pisos térmicos presentados en este anexo se determinaron de acuerdo con el Anexo No.\ 2 de la Resolución 194 de 2025 del Ministerio de Vivienda, Ciudad y Territorio, que define los parámetros y lineamientos de construcción sostenible y adopta los rangos climáticos municipales como base para el diseño de edificaciones eficientes energéticamente \parencite{UPME2025}. Para efectos del análisis desarrollado en el Capítulo~\ref{chap:diagnostico}, se tomaron los municipios de los siete departamentos de la región Caribe donde Air-e y Afinia tienen reportes de suscriptores en el SUI, filtrando únicamente los pisos térmicos cálido húmedo y cálido seco, que corresponden al 92\% de los municipios del área de análisis y a los contextos climáticos más representativos de los BS caracterizados en esta monografía. El listado completo de los 195 municipios con su clasificación climática se encuentra disponible en el archivo complementario de datos de esta monografía.


% ══════════════════════════════════════════════════════════════
\chapter{Datos de Consumo Residencial SUI -- Plataforma O3 (2023--2025)}
\label{anx:consumo_sspd}
% ══════════════════════════════════════════════════════════════

Los datos de consumo residencial presentados en la Sección~\ref{subsec:consumo_real} del Capítulo~\ref{chap:diagnostico} fueron obtenidos de la plataforma O3 dispuesta por la SSPD, filtrando la información bajo los criterios descritos en esa sección \parencite{SSPD2025}. La base de datos completa de consumos mensuales por municipio, departamento y piso térmico para el periodo 2023--2025 se encuentra disponible en el archivo complementario de datos de esta monografía, al que se puede acceder mediante el enlace indicado en el documento principal. Los valores promedio reportados en el texto —235, 232 y 228~kWh-mes para clima cálido húmedo, y 218, 214 y 211~kWh-mes para clima cálido seco— corresponden al promedio ponderado por número de suscriptores de los municipios incluidos en el filtro.